\documentclass[conference,letterpaper]{IEEEtran}
\IEEEoverridecommandlockouts
% Toque MCIC sobre IEEEtran: solo color y diagramas. No cambia márgenes,
% fuentes, columnas ni tamaños de la plantilla IEEE.
% tectonic compila con XeTeX: Times New Roman (la de IEEE y la secundaria del manual UD)
% se lee de las fuentes del manual de marca; si no está, TeX Gyre Termes (clon de Times).
\usepackage{amsmath,amssymb}
\usepackage{newtxmath}
\usepackage{fontspec}
\providecommand{\UDmarca}{/Users/juferoga/repos/ud/maestria/mcic/.UniversidadDistrital-manual-marca}
\IfFileExists{\UDmarca/fuentes/times.ttf}{%
  \setmainfont{times.ttf}[Path=\UDmarca/fuentes/, BoldFont=timesbd.ttf,
    ItalicFont=timesi.ttf, BoldItalicFont=timesbi.ttf]}{%
  \setmainfont{texgyretermes}[Extension=.otf, UprightFont=*-regular, BoldFont=*-bold,
    ItalicFont=*-italic, BoldItalicFont=*-bolditalic]}
\setmonofont{texgyrecursor}[Extension=.otf, UprightFont=*-regular, BoldFont=*-bold, Scale=0.95]
\usepackage[spanish,es-noshorthands,es-tabla,es-nolists]{babel}
\usepackage{graphicx}
\usepackage{booktabs,tabularx,array,multirow}
\usepackage[table]{xcolor}
\usepackage{tikz}
\usetikzlibrary{positioning,arrows.meta,calc,fit}
\usepackage[hidelinks]{hyperref}
\usepackage{cite}
\usepackage{url}

% Paleta del manual de marca UD (design.md)
\definecolor{udazul}{HTML}{1291C0}
\definecolor{udazulosc}{HTML}{0C6F96}
\definecolor{uddorado}{HTML}{B49739}
\definecolor{udverde}{HTML}{00A859}
\definecolor{udplata}{HTML}{98999A}
\definecolor{udsable}{HTML}{373435}
\definecolor{udrojo}{HTML}{ED3237}

\hypersetup{colorlinks=true,linkcolor=udazulosc,citecolor=udazulosc,urlcolor=udazulosc}

% Bloques de la ecuación de búsqueda: el mismo color en ecuación, PRISMA y matriz.
\newcommand{\bloqueA}[1]{\textcolor{udazulosc}{#1}}
\newcommand{\bloqueB}[1]{\textcolor{uddorado!85!black}{#1}}
\newcommand{\bloqueC}[1]{\textcolor{udverde!80!black}{#1}}
\newcommand{\etiqueta}[2]{\tikz[baseline=(t.base)]\node[fill=#1,text=white,inner sep=1.6pt,
  rounded corners=1pt,font=\bfseries\scriptsize](t){#2};}

% Cabecera de tablas: filete azul y fila sombreada clara
\newcommand{\cabecera}{\rowcolor{udazul!12}}
\arrayrulecolor{udazulosc}

% Columnas de la matriz
\newcolumntype{L}[1]{>{\raggedright\arraybackslash}p{#1}}
\newcolumntype{Y}{>{\raggedright\arraybackslash}X}

\makeatletter
\providecommand{\filas}[1]{\@@input #1 }
\makeatother

% Marca de enfoque en la matriz (A..E)
\newcommand{\enf}[1]{\textcolor{udazulosc}{\bfseries #1}}

% Cajas del PRISMA
\tikzset{
  caja/.style={draw=udazulosc,line width=0.5pt,rounded corners=1.5pt,align=center,
    font=\scriptsize,inner sep=2.5pt,text width=#1},
  caja/.default=3.6cm,
  excl/.style={draw=udplata,line width=0.5pt,rounded corners=1.5pt,align=left,
    font=\scriptsize,inner sep=2.5pt,fill=udplata!8,text width=#1},
  excl/.default=2.75cm,
  fase/.style={fill=udazul,text=white,font=\bfseries\tiny,minimum width=0.42cm,
    rounded corners=1pt,inner sep=1pt},
  flecha/.style={-{Latex[length=1.6mm]},udsable,line width=0.5pt},
}

% Generado por revision/generar_e1.py; no editar a mano.
\newcommand{\nArxivQuno}{35}
\newcommand{\nArxivQdos}{3}
\newcommand{\nArxivTotal}{38}
\newcommand{\nArxivDup}{0}
\newcommand{\nOaQuno}{162}
\newcommand{\nOaQdos}{33}
\newcommand{\nOaTotal}{195}
\newcommand{\nOaDup}{39}
\newcommand{\nDup}{39}
\newcommand{\nTamizados}{194}
\newcommand{\nExclTitulo}{147}
\newcommand{\nPasan}{47}
\newcommand{\nExclTexto}{10}
\newcommand{\nIncluidosBases}{37}
\newcommand{\nOtros}{21}
\newcommand{\nTotalIncluidos}{58}
\newcommand{\nMatriz}{16}
\newcommand{\nMatrizRecientes}{13}
\newcommand{\nRecientes}{49}
\newcommand{\nTextoCompleto}{35}
\newcommand{\nSoloResumen}{23}
\newcommand{\fechaConsulta}{2026-11-06}
\newcommand{\fechaConsultaOa}{2026-11-07}

\renewcommand{\IEEEkeywordsname}{Palabras clave}
\newcommand{\pendiente}[1]{\textcolor{udrojo}{\textbf{[#1]}}}
\newcommand{\emojiSmile}{\tikz[baseline=-0.6ex]{%
  \draw[fill=yellow!95!orange,draw=orange!90!black,line width=0.5pt] (0,0) circle (1.05ex);%
  \fill[black] (-0.35ex,0.23ex) circle (0.14ex);%
  \fill[black] (0.35ex,0.23ex) circle (0.14ex);%
  \draw[line width=0.55pt,black,line cap=round] (-0.48ex,-0.11ex) arc[start angle=200,end angle=340,radius=0.51ex];%
}}
\raggedbottom

\begin{document}

\title{Riesgo de los commits del kernel de Linux bajo deriva temporal: propuesta de un clasificador difuso-evolutivo interpretable y estado del arte}

\author{\IEEEauthorblockN{Juan Felipe Rodríguez Galindo}
\IEEEauthorblockA{Maestría en Ciencias de la Información y las Comunicaciones\\
Universidad Distrital Francisco José de Caldas, Bogotá, Colombia\\
Código 20261595004 \textperiodcentered{} Inteligencia Computacional \textperiodcentered{} Entregable 1}}

\maketitle

\begin{abstract}
Predecir si un commit del kernel de Linux requerirá corrección permite priorizar la revisión de código aprovechando el trailer \texttt{Fixes:}. Este trabajo formaliza el problema, la pregunta y el criterio de éxito del proyecto, y aporta una revisión estructurada tipo PRISMA (\nTotalIncluidos{} estudios incluidos, 31 citados aquí; la matriz completa va en el anexo). Se evidencia la falta de clasificadores de reglas difusas optimizados evolutivamente a nivel de commit bajo validación temporal y evaluados con AUC-PR, así como el estudio de su degradación ante la deriva en 2026.
\end{abstract}

\begingroup\renewcommand{\abstractname}{Abstract}
\begin{abstract}
Predicting Linux kernel commit risk via \texttt{Fixes:} trailers focuses code review. We define the problem, question, and success criteria, backed by a PRISMA review (\nTotalIncluidos{} included studies, 31 cited here; the full matrix is in the appendix). No prior study evaluates an evolutionarily tuned fuzzy rule classifier at commit level using temporal validation, AUC-PR, and drift analysis in 2026.
\end{abstract}\endgroup

\begin{IEEEkeywords}
predicción de defectos justo a tiempo, kernel de Linux, deriva temporal, sistemas difusos, computación evolutiva, AUC-PR.
\end{IEEEkeywords}

%=====================================================================
\section{Introducción}

Desde hace dos décadas, el riesgo de un commit se estima con métricas del propio cambio: cuánto toca, cuánto se difunde y la experiencia del autor \cite{Mockus2000,Kamei2013}. Es la predicción justo a tiempo (JIT): la alerta llega cuando el cambio aún está en revisión. El kernel de Linux cuenta con una ventaja clave: cuando un commit corrige un fallo previo, el trailer \texttt{Fixes:} identifica explícitamente al culpable sin recurrir a heurísticas como SZZ mediante \texttt{git blame} \cite{Sliwerski2005}. En estudios empíricos, el 81\,\% se confirmó como causa del bug; el resto es ambiguo o incorrecto \cite{Lyu2024}.

En 2026 cambió quién escribe parte de estos cambios. La versión 7.0 pidió declarar con \texttt{Assisted-by:} la ayuda de asistentes de IA \cite{KernelAdd}, simplificada luego a \texttt{Assisted-by: LLM} \cite{KernelSimpl}. El 14 de octubre se retiró al constatar que solo cerca del 10\,\% de los commits la incluía, pese a un uso real mucho mayor \cite{KernelRemove}. Por ello, la <<era de la IA>> se define por frontera temporal (Fig.~\ref{fig:tiempo}) y no por la etiqueta, que queda solo como referencia complementaria.

\begin{figure}[!t]
\centering
% Fig. 1: ventanas de datos frente a los hitos del trailer Assisted-by.
% Eje izquierdo esquemático (2019-2025); 2026 a escala por meses (0,58 cm por mes).
\begin{tikzpicture}[x=1cm,y=1cm,font=\scriptsize]
  \def\m{0.5}           % ancho de un mes de 2026
  \def\xo{2.6}          % inicio de 2026 en el eje
  % --- bandas de datos
  \fill[udazul!18] (0,0) rectangle (\xo,0.42);
  \node[udazulosc,align=center,font=\tiny\bfseries] at (\xo/2,0.21) {entrenamiento y validación\\ móvil (2019--2025)};
  \fill[uddorado!30] (\xo,0) rectangle ({\xo+7*\m},0.42);
  \node[uddorado!60!black,font=\tiny\bfseries] at ({\xo+3.5*\m},0.21) {prueba: ene--jul 2026};
  \fill[udplata!25] ({\xo+7*\m},0) rectangle ({\xo+10.2*\m},0.42);
  \node[udsable,align=center,font=\tiny] at ({\xo+8.6*\m},0.21) {censura\\ ($<$90 d)};
  % --- eje
  \draw[udsable,line width=0.5pt] (0,0) -- ({\xo+10.4*\m},0);
  \draw[udsable,line width=0.5pt] (\xo-0.42,-0.08) -- (\xo-0.34,0.08);
  \draw[udsable,line width=0.5pt] (\xo-0.34,-0.08) -- (\xo-0.26,0.08);
  \foreach \x/\t in {0/2019, 1.9/2025} \node[below,udsable] at (\x,0) {\t};
  \foreach \i/\t in {0/ene,2/mar,4/may,6/jul,8/sep,10/nov}
    {\draw[udsable] ({\xo+\i*\m},0) -- ++(0,-0.06);
     \node[below,udsable] at ({\xo+\i*\m},-0.02) {\t};}
  \node[below=7pt,udsable,font=\tiny\bfseries] at ({\xo+5*\m},0) {2026};
  % --- hitos del trailer (encima del eje)
  \foreach \x/\t/\yy in {0.17/v7.0: se pide \texttt{Assisted-by}/0.95,
                          7.07/v7.3: solo ``\texttt{LLM}''/1.35,
                          9.45/v7.4: se elimina/0.95}{
    \draw[udrojo,line width=0.6pt] ({\xo+\x*\m},0.42) -- ({\xo+\x*\m},\yy-0.08);
    \fill[udrojo] ({\xo+\x*\m},\yy-0.08) circle (1.1pt);
    \node[udrojo,above,font=\tiny,inner sep=1pt] at ({\xo+\x*\m},\yy-0.06) {\t};}
  % --- corte de observación
  \draw[udsable,dashed] ({\xo+10.2*\m},-0.05) -- ({\xo+10.2*\m},0.6)
     node[above,font=\tiny,udsable,inner sep=1pt] {corte};
\end{tikzpicture}

\caption{Diseño temporal propuesto frente a los hitos de \texttt{Assisted-by} \cite{KernelAdd,KernelSimpl,KernelRemove}. Los commits posteriores a julio de 2026 quedan censurados: no han tenido los 90 días para recibir un \texttt{Fixes:}.}
\label{fig:tiempo}
\end{figure}

\subsection{Problema y pregunta de investigación}

\textbf{Problema.} Clasificar cada commit del kernel como \emph{riesgoso} si un commit posterior lo cita con \texttt{Fixes:} en 90 días o menos, usando solo información disponible en el momento de integrarlo, y con un modelo que un mantenedor pueda leer.

\textbf{Pregunta.} ¿Un clasificador de a lo sumo 15 reglas difusas, cuyas funciones de pertenencia y reglas se ajustan con un algoritmo evolutivo, alcanza en los commits de 2026 un AUC-PR de al menos el 90\,\% del mejor modelo de referencia (regresión logística o \emph{gradient boosting}), y pierde menos AUC-PR que ese modelo al pasar de la validación (ventanas anuales 2023, 2024 y 2025, entrenando desde 2019) a la prueba (enero--julio de 2026)?

\subsection{Criterio de éxito (etapa 1 de CRISP-ML(Q) \cite{Studer2021})}

Sea $\mathrm{AP}_m^{V}$ el AUC-PR medio del modelo $m$ en las ventanas de validación y $\mathrm{AP}_m^{T}$ su AUC-PR en la prueba de 2026, ambos con intervalo de confianza del 95\,\% por \emph{bootstrap}. Sean $\Delta_m=(\mathrm{AP}_m^{V}-\mathrm{AP}_m^{T})/\mathrm{AP}_m^{V}$ la caída relativa, $\pi^{T}$ la prevalencia de positivos en la prueba (el AUC-PR de la clase mayoritaria), FE el modelo propuesto y $R$ la mejor referencia clásica en validación. El proyecto se declara exitoso si se cumplen las tres condiciones de la Tabla~\ref{tab:criterio}. Si no se cumplen, se reporta así, sin volver a ajustar nada con la prueba.

\begin{table}[!t]
\caption{Criterio de éxito, fijado antes de ver los datos}
\label{tab:criterio}
\centering\footnotesize
\begin{tabularx}{\columnwidth}{@{}l Y@{}}
\toprule
\cabecera & \textbf{Condición}\\
\midrule
C1 Desempeño & $\mathrm{AP}_{FE}^{T}\ge 0{,}90\,\mathrm{AP}_{R}^{T}$\\
C2 Robustez & $\Delta_{FE}<\Delta_{R}$ (diferencia observada, con IC 95\,\%)\\
C3 Piso & El IC 95\,\% de $\mathrm{AP}_{FE}^{T}$ excluye $\pi^{T}$; McNemar contra cada referencia con $\alpha=0{,}05$\\
Restricción & $\le$15 reglas y $\le$3 antecedentes por regla\\
\bottomrule
\end{tabularx}
\end{table}

\subsection{Datos y consideraciones éticas}
Los datos son públicos: el historial oficial de \texttt{torvalds/linux} (GPL-2.0), minado directamente para un registro continuo de 2014 a 2026. Para contrastar la etiqueta se recurre además al dataset de Maes Bermejo \emph{et al.} en Zenodo \cite{MaesBermejo2024}. Por ser fuentes enteramente abiertas no se requiere aprobación de un comité de ética; por privacidad, los nombres y correos de los autores se sustituyen por identificadores anónimos.


%=====================================================================
\section{Estado del arte}

\subsection{Método de búsqueda}
La búsqueda sigue las directrices de PRISMA 2020 \cite{Page2021} en versión sintetizada. El problema se descompone en tres componentes temáticos esenciales, unidos con OR por dentro y AND entre ellos:

\noindent\fbox{\parbox{\dimexpr\columnwidth-2\fboxsep-2\fboxrule}{\footnotesize\raggedright
\etiqueta{udazulosc}{A} \bloqueA{(commit* OR ``code change*'' OR ``software change*'' OR patch* OR ``pull request*'' OR ``Linux kernel'')} AND\\
\etiqueta{uddorado}{B} \bloqueB{(``defect prediction'' OR ``bug prediction'' OR ``bug-introducing'' OR ``bug-inducing'' OR ``fix-inducing'' OR ``just-in-time'' OR ``change risk'' OR ``risk prediction'' OR ``defect-prone'' OR buggy)} AND\\
\etiqueta{udverde}{C} \bloqueC{(fuzzy OR ``rule-based'' OR evolutionary OR ``genetic algorithm'' OR ``genetic programming'' OR ``search-based'' OR interpretab* OR explainab*)}}}

\smallskip
La ecuación Q1 combina los tres bloques. La Q2 cambia el bloque C por deriva temporal (``concept drift'', drift, ``temporal validation'', ``time-aware'', ``realistic evaluation'', ``data shift'', ``distribution shift'', ``time travel''). Ambas consultas se ejecutaron sobre 2019--2026 vía APIs públicas en arXiv y OpenAlex (título y resumen) para garantizar reproducibilidad. OpenAlex indexa IEEE, ACM, Springer y Elsevier, supliendo la falta de acceso institucional a Scopus o IEEE Xplore. Clásicos previos a 2019 o hallados por citas entran como ``otras fuentes''.

Los criterios se definieron antes del tamizaje. Se incluye si: (I1) predice el riesgo de un cambio o estudia la validez de esa evaluación (orden temporal, latencia, ruido), (I2) reporta métricas cuantitativas, y (I3) describe sus datos. Se excluye si: (E1) trata generación o reparación de código, (E2) evalúa solo a nivel de archivo o módulo (se admitieron propuestas difusas y evolutivas de módulo ante la ausencia a nivel de commit), (E3) es de otro dominio, (E4) es solo un dataset sin modelo, o (E5) no se puede verificar. De las bases se evaluaron \nPasan{} registros en elegibilidad (\nExclTexto{} excluidos y \nIncluidosBases{} incluidos); a estos se sumaron \nOtros{} de otras fuentes bajo los mismos criterios I y E, totalizando \nTotalIncluidos{} estudios incluidos (37 citados aquí; la matriz completa va en el anexo). De ellos, \nTextoCompleto{} se leyeron a texto completo y \nSoloResumen{} por resumen o fuentes secundarias debido a restricciones de acceso institucional. El flujo se resume en la Fig.~\ref{fig:prisma}.

\begin{figure}[!t]
\centering
% Fig. 2: diagrama de flujo tipo PRISMA 2020. Todos los números son macros de generado/conteos.tex.
\begin{tikzpicture}[node distance=2.6mm and 2.4mm]
  % Identificación
  \node[caja=3.75cm] (bases) {\textbf{Bases de datos} (2019--2026)\\
     arXiv: Q1 \nArxivQuno{} $+$ Q2 \nArxivQdos{} $=$ \nArxivTotal\\
     OpenAlex: Q1 \nOaQuno{} $+$ Q2 \nOaQdos{} $=$ \nOaTotal};
  \node[caja=2.45cm,right=of bases] (otros) {\textbf{Otras fuentes}\\ clásicos, citas y
     búsqueda dirigida: \nOtros};
  % Tamizaje
  \node[caja=3.75cm,below=4mm of bases] (tam) {\textbf{Tamizaje por título y resumen}\\
     \nTamizados{} registros (sin \nDup{} duplicados)};
  \node[excl,right=of tam] (extam) {Excluidos (E1--E5): \nExclTitulo};
  \node[caja=3.75cm,below=of tam] (tc) {\textbf{Evaluados en elegibilidad}: \nPasan};
  \node[excl,right=of tc] (extc) {Excluidos tras leer: \nExclTexto};
  \node[caja=3.75cm,below=of tc,fill=udazul!10] (inc) {\textbf{Incluidos: \nTotalIncluidos}
     (\nIncluidosBases{} de bases $+$ \nOtros{} de otras fuentes)\\ \nRecientes{} desde 2021;
     \nMatriz{} en la Tabla~\ref{tab:matriz}};
  \draw[flecha] (bases) -- (tam);
  \draw[flecha] (tam) -- (tc);
  \draw[flecha] (tc) -- (inc);
  \draw[flecha] (otros.east) -- ++(1.6mm,0) |- (inc.east);
  \draw[flecha,udplata] (tam.east) -- (extam.west);
  \draw[flecha,udplata] (tc.east) -- (extc.west);
  % Etiquetas de fase
  \foreach \n/\t in {bases/Ident., tam/Tamiz., inc/Incl.}
    \node[fase,rotate=90,anchor=south] at ([xshift=-1.2mm]\n.west) {\t};
\end{tikzpicture}

\caption{Flujo tipo PRISMA 2020 de la búsqueda. Conteos generados desde los registros de la búsqueda.}
\label{fig:prisma}
\end{figure}

% Tabla II: matriz reducida del estado del arte. Cifras tomadas de revision/fichas/*.json
% (cada una con su URL verificada); la matriz completa está en revision/resultados/matriz_completa.csv.
\begin{table*}[!t]
\caption{Matriz del estado del arte (\nMatriz{} de los \nTotalIncluidos{} incluidos; la matriz completa va en el anexo). Enfoques: A métricas y aprendizaje clásico, B profundo y LLM, C interpretable, difuso o evolutivo, D evaluación realista y deriva, E etiqueta. $^\dagger$Solo nivel de módulo. $^\ddagger$Preprint.}
\label{tab:matriz}
\centering
\fontsize{5.6}{6.6}\selectfont
\setlength{\tabcolsep}{2pt}
\renewcommand{\arraystretch}{1.08}
\begin{tabularx}{\textwidth}{@{}L{1.05cm}L{2.45cm}L{2.3cm}L{1.75cm}L{1.75cm}L{2.55cm}L{1.95cm}Y@{}}
\toprule
\cabecera \textbf{Referencia} & \textbf{Problema y datos} & \textbf{Paradigma y modelo} & \textbf{Referencias comparadas} & \textbf{Validación} & \textbf{Métrica y resultado} & \textbf{Limitaciones} & \textbf{Relevancia para el proyecto}\\
\midrule
\enf{A} Kamei \cite{Kamei2013} (2013) & 11 proyectos (6 libres, 5 comerciales), $>$250 mil cambios; SZZ & Reg. logística, 14 métricas de cambio & Azar & VC aleatoria 10 pliegues & Precisión 34\,\%, sensib. 64\,\%; AUC 0,73 (comerciales) & VC mezcla pasado y futuro & Variables base y referencia logística\\
\enf{B} Zeng \cite{Zeng2021} (2021) & 310 mil commits, 6 proyectos (Qt, OpenStack, Go\ldots) & Réplica de DeepJIT y CC2Vec; LApredict (RL con una variable) & DeepJIT, CC2Vec, DBN-JIT, LR-JIT & Cronológica 80/20; recorta la cola & AUC media: LApredict 0,725; DeepJIT 0,707 & Una corrida; solo AUC & Lo simple iguala a lo profundo con partición temporal\\
\enf{B} Ni \cite{Ni2022} (2022) & JIT-Defects4J: 27\,319 commits Java, 8,5\,\% positivos & JIT-Fine: CodeBERT + 14 métricas & LApredict, DeepJIT, CC2Vec, JITLine & Temporal 80/20 por proyecto & F1 0,431; AUC 0,881 (JITLine 0,261/0,802) & Solo Java; costo no reportado & Cota alta de AUC; modelo opaco\\
\enf{B} Liu \cite{Liu2026} (2026) & Diffs de Meta; cambios que causaron caídas ($\approx$2 por mil) & LLM (Diff Risk Score) + atención por bloque & Ninguna cuantitativa & No descrita & Top-2 bloques cubren 53,9\,\% de líneas causantes & Una empresa; 48 de 55 opiniones negativas & Explicación post hoc frente a reglas intrínsecas\\
\enf{C} JITLine \cite{JITLine2021} (2021) & Qt y OpenStack, 37\,524 commits & RF + tokens; SMOTE ajustado por evolución diferencial; LIME & EALR, DeepJIT, CC2Vec & Una partición temporal & AUC 0,83/0,82; F1 0,33/0,24 & Una partición, sin pruebas estadísticas & Usa evolución y explica, pero sobre caja negra\\
\enf{C} PyExplainer \cite{PyExplainer2021} (2021) & Qt y OpenStack, 40\,798 commits & Reglas locales (RuleFit) sobre RF y RL & LIME & Temporal 70/30 & RF: AUC 0,75/0,74; reglas 18--38\,\% más fieles que LIME & Dos proyectos; vecinos aleatorios & Reglas post hoc, no un modelo de reglas\\
\enf{C} Aftab \cite{Aftab2023} (2023)$^\dagger$ & NASA, 5 conjuntos, 3\,579 módulos & Inferencia difusa (8 reglas) que combina NB, RNA y AD & NB, RNA, AD & Aleatoria 70/30 & Exactitud 91,05\,\%; sensib. 0,33 & Reglas a mano; exactitud con desbalance & Muestra por qué no usar exactitud\\
\enf{C} Dey \cite{Dey2025} (2025)$^\dagger$ & NASA: JM1, KC1, PC1, KC2, CM1 & Selección multiobjetivo (estrella de mar) + integral difusa de Choquet & 7 clasificadores; MOGA, MOGWO & VC aleatoria 10 pliegues & Exactitud 0,85--0,96; sin AUC & Nivel módulo; F1 y sensib. casi perfectas con 7--20\,\% positivos & Lo difuso-evolutivo reciente no llega al commit\\
\enf{D} McIntosh \cite{McIntosh2018} (2018) & Qt y OpenStack, 37\,524 cambios (8 y 13\,\%) & RL no lineal con splines & Periodo corto frente a largo & Longitudinal en el tiempo & AUC cae 11--34 puntos en un año; reentrenar recupera 6--24 & Dos sistemas; filtros de SZZ & Evidencia de deriva; diseño a replicar\\
\enf{D} Cabral \cite{Cabral2019} (2019) & 10 proyectos GitHub, 2003--2018 & En línea: ORB (árboles de Hoeffding) & OOB, UOB, ventana deslizante & Precuencial con espera de 90 días & G-mean de JIT fijo hasta 97,2\,\% menor que OOB & G-mean, sin AUC-PR & Sustenta la ventana de 90 días\\
\enf{D} Olewicki \cite{Olewicki2024} (2024) & Ubisoft: 4 proyectos de riesgo, 50--82 mil commits & MLP con aprendizaje continuo & XGBoost reentrenado, ingenuo & Secuencial en el tiempo & F1 +9 a +31\,\% en 6 de 8 proyectos & Una empresa; sin latencia & Deriva industrial en riesgo de commits\\
\enf{D} Nguyen \cite{Nguyen2025} (2025) & Linux (993\,704 commits) y FFmpeg; Linux con \texttt{Fixes:} & 8 técnicas JIT de vulnerabilidades (RL, DeepJIT, JIT-Fine\ldots) & Escenario idealizado frente a realista & Temporal 75/5/20 & PR-AUC en Linux 0,805 $\rightarrow$ 0,016 ($-$98\,\%) & Vulnerabilidades, no defectos & El más cercano: Linux, \texttt{Fixes:}, AUC-PR\\
\enf{D} Mazloomzadeh \cite{Mazloomzadeh2026} (2026)$^\ddagger$ & AIDev: 489 repos. Python, 9\,428 PR de agentes & Estudio longitudinal (GLMM, SZZ) & PR humanos & Por trimestres & Sin diferencia relevante en inducción de bugs ($V=0{,}068$) & SZZ con P=0,39 & La era IA cambia las covariables\\
\enf{E} Lyu \cite{Lyu2024} (2024) & KFC: 76\,046 commits con \texttt{Fixes:}, Linux 2013--2022 & 6 variantes de SZZ + TC-SZZ & Rosa et al. & Contra \texttt{Fixes:} como verdad & R-SZZ F1 0,51; 81\,\% de \texttt{Fixes:} confirmados & \texttt{Fixes:} es opcional & Valida la etiqueta del proyecto\\
\enf{E} Herbold \cite{Herbold2022} (2022) & 398 versiones, 38 proyectos Apache & Calidad de datos; NB entre proyectos & SZZ a 6 meses frente a etiqueta validada & Entre proyectos; Friedman + Nemenyi & Cerca de la mitad de las correcciones de SZZ no lo son & Solo Apache & Ruido de SZZ; preferir \texttt{Fixes:}\\
\enf{E} Falessi \cite{Falessi2022} (2022) & 19 proyectos Apache, $>$4\,000 defectos & 15 clasificadores de Weka & Con y sin defectos latentes & Por versiones, 66/33 & Pérdida mediana de sensib. $\approx$90\,\% por defectos latentes & Nivel de clase & Censura: recortar la cola reciente\\
\bottomrule
\end{tabularx}
\end{table*}


\subsection{Síntesis por enfoques}

\textbf{A. Métricas de cambio y aprendizaje clásico.} La línea fundacional usa métricas del commit con regresión logística \cite{Mockus2000,Kamei2013}. Sigue siendo un punto de referencia muy sólido: con partición cronológica estricta, una regresión logística de una sola variable (líneas añadidas) iguala o supera en AUC a arquitecturas profundas como DeepJIT y CC2Vec \cite{Zeng2021}.

\textbf{B. Aprendizaje profundo y LLM.} Redes sobre el diff y modelos preentrenados con métricas \cite{Ni2022} alcanzan techos de AUC-ROC cercanos a 0,88 en JIT-Defects4J. Tienen tres limitaciones clave: son cajas negras opacas, se evalúan casi exclusivamente en Java y rara vez reportan AUC-PR. En la industria, aunque se han usado LLMs para puntuar diffs, sus mapas de atención generaron baja confianza en los desarrolladores \cite{Liu2026}.

\textbf{C. Interpretabilidad, difuso y evolutivo.} Los profesionales piden explicaciones claras para cada alerta \cite{Jiarpakdee2021}. La respuesta habitual son explicadores \emph{post hoc}: LIME en JITLine \cite{JITLine2021}, reglas locales en PyExplainer \cite{PyExplainer2021} o commits afines en IRJIT \cite{Sahar2024}, los cuales carecen de garantías formales de consistencia \cite{Yu2024FoX}. A nivel de commit, lo evolutivo y difuso se reducen a tareas accesorias: evolución diferencial para SMOTE \cite{JITLine2021}, NSGA-II para calibrar parámetros logísticos \cite{Chen2018MULTI}, optimización de hiperparámetros \cite{Kang2021,Mostafa2025} o Fuzzy C-Means para modular capas de atención \cite{Chen2025}. A nivel de módulo existen modelos neuro-difusos sobre datos NASA con particiones aleatorias y exactitud como métrica \cite{Aftab2023,Dey2025} (donde una exactitud de 91\,\% enmascara una sensibilidad de apenas 0,33 \cite{Aftab2023}). No se identificó ningún clasificador de reglas difusas para riesgo de commits.

\textbf{D. Evaluación realista, deriva y latencia.} Los modelos JIT se degradan sustancialmente en el tiempo. En Qt y OpenStack, el AUC cae entre 11 y 34 puntos en solo un año sin reentrenar \cite{McIntosh2018}. La validación cruzada infla los resultados: JITBoost cae de 0,891 a 0,863 de AUC al evaluar en partición temporal \cite{Shehab2023}. La deriva en JIT se aborda mediante detección continua \cite{Chitsazian2026}, reciclaje de modelos \cite{Patel2024}, actualización en línea continua \cite{Olewicki2024} o ventanas de espera \cite{Tian2026}. Sin embargo, casi todos reportan AUC-ROC y omiten AUC-PR. La etiqueta llega tarde (latencia de 15 a más de 400 días \cite{Cabral2019}), lo que distorsiona las métricas hasta un 25\,\% si no se depura la cola temporal \cite{Song2023}. En el kernel de Linux este efecto es drástico: al pasar de datos balanceados artificialmente a la prevalencia real, el AUC-PR de ocho métodos cae de 0,805 a 0,016 \cite{Nguyen2025}. Además, los commits de agentes no elevan la tasa de defectos, pero alteran el tamaño y estructura de los cambios \cite{Mazloomzadeh2026}, ubicando la deriva en las covariables.

\textbf{E. Fiabilidad de la etiqueta.} SZZ sufre de alto ruido: casi la mitad de los commits que señala como correctores no lo son \cite{Herbold2022}, y el filtrado con LLM reduce 21\,\% los casos sospechosos en PostgreSQL \cite{Veras2025}. En Linux, variantes de SZZ alcanzan un F1 de apenas 0,49--0,51 frente a la verdad del trailer \texttt{Fixes:} \cite{Lyu2024}. Defectos latentes aún no descubiertos bajan la sensibilidad cerca de un 90\,\% si no se aplica una ventana prudencial de censura temporal \cite{Falessi2022}.

\subsection{Brecha}
Los enfoques anteriores dejan cuatro vacíos concretos:
\begin{enumerate}
\item Inexistencia de un clasificador de reglas difusas optimizado mediante algoritmos evolutivos para predecir riesgo a nivel de commit.
\item Escasa medición de la deriva temporal con AUC-PR y prevalencia real en commits del kernel respaldados por trailers \texttt{Fixes:}.
\item Falta de estudios sobre si la interpretabilidad intrínseca penaliza el desempeño o aporta robustez frente a la deriva.
\item 2026 es el primer año en que el kernel pidió declarar la ayuda de IA \cite{KernelAdd}, aunque la etiqueta se retiró en octubre por su bajo uso \cite{KernelRemove}.
\end{enumerate}


%=====================================================================
\section*{Anexos}
La matriz completa de los 58 trabajos incluidos se adjunta en Classroom como archivo de Excel. Mientras creo el repositorio con el código y los datos, los anexos se entregan por ese medio.

%=====================================================================
\section*{Declaración de uso de IA generativa}
{\footnotesize Usé Gemini para revisar y corregir errores y para redactar este informe a partir de los resultados del cuaderno y la recuperación de algunos papers para la evaluación del mismo.}


% Generado por revision/generar_e1.py 1 en orden de cita; no editar a mano.
\begin{thebibliography}{37}
\bibitem{Mockus2000} A. Mockus and D. M. Weiss, ``Predicting risk of software changes,'' Bell Labs Tech. J., vol. 5, no. 2, pp. 169--180, Apr.--Jun. 2000, doi: \mbox{10.1002/bltj.2229}.
\bibitem{Kamei2013} Y. Kamei, E. Shihab, B. Adams, A. E. Hassan, A. Mockus, A. Sinha, and N. Ubayashi, ``A large-scale empirical study of just-in-time quality assurance,'' IEEE Trans. Softw. Eng., vol. 39, no. 6, pp. 757--773, Jun. 2013, doi: \mbox{10.1109/TSE.2012.70}.
\bibitem{Sliwerski2005} J. Śliwerski, T. Zimmermann, and A. Zeller, ``When do changes induce fixes?'' in Proc. Int. Workshop Mining Softw. Repositories (MSR '05), St. Louis, MO, USA, 2005, pp. 1--5, doi: \mbox{10.1145/1083142.1083147}.
\bibitem{Lyu2024} Y. Lyu, H. J. Kang, R. Widyasari, J. Lawall, and D. Lo, ``Evaluating SZZ implementations: An empirical study on the Linux kernel,'' IEEE Trans. Softw. Eng., vol. 50, no. 9, pp. 2219--2239, Sep. 2024, doi: \mbox{10.1109/TSE.2024.3406718}.
\bibitem{KernelAdd} S. Levin, ``docs: add AI Coding Assistants documentation,'' Linux kernel, commit 78d979db6cef, Jan. 6, 2026. [Online]. Available: \url{https://github.com/torvalds/linux/commit/78d979db6cef}
\bibitem{KernelSimpl} C. Brauner, ``coding-assistants: simplify attribution,'' Linux kernel, commit 816d9992d9ed, Aug. 3, 2026. [Online]. Available: \url{https://github.com/torvalds/linux/commit/816d9992d9ed}
\bibitem{KernelRemove} S. Levin, ``docs: remove Assisted-by references,'' Linux kernel, commit 1e035297ffda, Oct. 14, 2026. [Online]. Available: \url{https://github.com/torvalds/linux/commit/1e035297ffda}
\bibitem{Studer2021} S. Studer, T. B. Bui, C. Drescher, A. Hanuschkin, L. Winkler, S. Peters, and K.-R. M\"uller, ``Towards CRISP-ML(Q): A machine learning process model with quality assurance methodology,'' Mach. Learn. Knowl. Extr., vol. 3, no. 2, pp. 392--413, 2021, doi: 10.3390/make3020020.
\bibitem{MaesBermejo2024} M. Maes Bermejo, J. M. González-Barahona, M. Gallego, and G. Robles, ``A dataset of Linux kernel commits'' (v3), Zenodo, 2024, doi: 10.5281/zenodo.10654193.
\bibitem{Page2021} M. J. Page \emph{et al.}, ``The PRISMA 2020 statement: An updated guideline for reporting systematic reviews,'' BMJ, vol. 372, Art. no. n71, 2021, doi: 10.1136/bmj.n71.
\bibitem{Zeng2021} Z. Zeng, Y. Zhang, H. Zhang, and L. Zhang, ``Deep just-in-time defect prediction: How far are we?'' in Proc. 30th ACM SIGSOFT Int. Symp. Softw. Testing Anal. (ISSTA), Virtual, Denmark, 2021, pp. 427--438, doi: \mbox{10.1145/3460319.3464819}.
\bibitem{Ni2022} C. Ni, W. Wang, K. Yang, X. Xia, K. Liu, and D. Lo, ``The best of both worlds: Integrating semantic features with expert features for defect prediction and localization,'' in Proc. 30th ACM Joint Eur. Softw. Eng. Conf. Symp. Found. Softw. Eng. (ESEC/FSE), Singapore, 2022, pp. 672--683, doi: \mbox{10.1145/3540250.3549165}.
\bibitem{Liu2026} Y. Liu, K. Jabre, R. Abreu, Z. J. Carmichael, V. Murali, A. Patel, J. Ge, W. Sun, C. Zhang, A. Mockus, D. Khavari, P. C. Rigby, and N. Nagappan, ``A preliminary study on explaining risk of code changes using LLM-based prediction models,'' in Proc. 3rd ACM Int. Conf. AI-Powered Softw. (AIware '26), Montreal, QC, Canada, 2026, pp. 259--267, doi: \mbox{10.1145/3805760.3814916}.
\bibitem{JITLine2021} C. Pornprasit and C. K. Tantithamthavorn, ``JITLine: A simpler, better, faster, finer-grained just-in-time defect prediction,'' in Proc. IEEE/ACM 18th Int. Conf. Mining Softw. Repositories (MSR), 2021, pp. 369--379, doi: \mbox{10.1109/MSR52588.2021.00049}.
\bibitem{PyExplainer2021} C. Pornprasit, C. Tantithamthavorn, J. Jiarpakdee, M. Fu, and P. Thongtanunam, ``PyExplainer: Explaining the predictions of just-in-time defect models,'' in Proc. 36th IEEE/ACM Int. Conf. Automated Softw. Eng. (ASE), 2021, pp. 407--418, doi: \mbox{10.1109/ASE51524.2021.9678763}.
\bibitem{Aftab2023} S. Aftab, S. Abbas, T. M. Ghazal, M. Ahmad, H. A. Hamadi, C. Y. Yeun, and M. A. Khan, ``A cloud-based software defect prediction system using data and decision-level machine learning fusion,'' Mathematics, vol. 11, no. 3, Art. no. 632, 2023, doi: \mbox{10.3390/math11030632}.
\bibitem{Dey2025} R. Dey, J. Piri, B. Acharya, P. P. Das, V. C. Gerogiannis, and A. Kanavos, ``A hybrid evolutionary fuzzy ensemble approach for accurate software defect prediction,'' Mathematics, vol. 13, no. 7, Art. no. 1140, 2025, doi: \mbox{10.3390/math13071140}.
\bibitem{McIntosh2018} S. McIntosh and Y. Kamei, ``Are fix-inducing changes a moving target? A longitudinal case study of just-in-time defect prediction,'' IEEE Trans. Softw. Eng., vol. 44, no. 5, pp. 412--428, May 2018, doi: \mbox{10.1109/TSE.2017.2693980}.
\bibitem{Cabral2019} G. G. Cabral, L. L. Minku, E. Shihab, and S. Mujahid, ``Class imbalance evolution and verification latency in just-in-time software defect prediction,'' in Proc. IEEE/ACM 41st Int. Conf. Softw. Eng. (ICSE), Montreal, QC, Canada, 2019, pp. 666--676, doi: \mbox{10.1109/ICSE.2019.00076}.
\bibitem{Olewicki2024} D. Olewicki, S. Habchi, M. Nayrolles, M. Faramarzi, S. Chandar, and B. Adams, ``On the costs and benefits of adopting lifelong learning for software analytics -- Empirical study on brown build and risk prediction,'' in Proc. 46th Int. Conf. Softw. Eng.: Softw. Eng. Pract. (ICSE-SEIP), Lisbon, Portugal, 2024, pp. 275--286, doi: \mbox{10.1145/3639477.3639717}.
\bibitem{Nguyen2025} D. Nguyen, T. Le-Cong, T. H. M. Le, M. A. Babar, and Q.-T. Huynh, ``Toward realistic evaluations of just-in-time vulnerability prediction,'' in Proc. IEEE Int. Conf. Softw. Maint. Evol. (ICSME), Auckland, New Zealand, 2025, pp. 1--10, doi: \mbox{10.1109/ICSME64153.2025.00023}.
\bibitem{Mazloomzadeh2026} I. Mazloomzadeh, M. M. Morovati, and F. Khomh, ``How do AI coding agents contribute to software development? An empirical study of agentic pull requests,'' arXiv:2607.21832, Jul. 2026, doi: \mbox{10.48550/arXiv.2607.21832}.
\bibitem{Herbold2022} S. Herbold, A. Trautsch, F. Trautsch, and B. Ledel, ``Problems with SZZ and features: An empirical study of the state of practice of defect prediction data collection,'' Empir. Softw. Eng., vol. 27, no. 2, Art. no. 42, 2022, doi: \mbox{10.1007/s10664-021-10092-4}.
\bibitem{Falessi2022} D. Falessi, A. Ahluwalia, and M. Di Penta, ``The impact of dormant defects on defect prediction: A study of 19 Apache projects,'' ACM Trans. Softw. Eng. Methodol., vol. 31, no. 1, pp. 1--26, Jan. 2022, doi: \mbox{10.1145/3467895}.
\bibitem{Jiarpakdee2021} J. Jiarpakdee, C. K. Tantithamthavorn, and J. Grundy, ``Practitioners' perceptions of the goals and visual explanations of defect prediction models,'' in Proc. IEEE/ACM 18th Int. Conf. Mining Softw. Repositories (MSR), 2021, pp. 432--443, doi: \mbox{10.1109/MSR52588.2021.00055}.
\bibitem{Sahar2024} H. Sahar, A. A. Bangash, A. Hindle, and D. Barbosa, ``IRJIT: A simple, online, information retrieval approach for just-in-time software defect prediction,'' Empir. Softw. Eng., vol. 29, no. 5, Art. no. 131, 2024, doi: \mbox{10.1007/s10664-024-10514-z}.
\bibitem{Yu2024FoX} J. Yu, M. Fu, A. Ignatiev, C. Tantithamthavorn, and P. Stuckey, ``A formal explainer for just-in-time defect predictions,'' ACM Trans. Softw. Eng. Methodol., vol. 33, no. 7, pp. 1--31, Sep. 2024, doi: \mbox{10.1145/3664809}.
\bibitem{Chen2018MULTI} X. Chen, Y. Zhao, Q. Wang, and Z. Yuan, ``MULTI: Multi-objective effort-aware just-in-time software defect prediction,'' Inf. Softw. Technol., vol. 93, pp. 1--13, Jan. 2018, doi: \mbox{10.1016/j.infsof.2017.08.004}.
\bibitem{Kang2021} J. Kang, S. Kwon, D. Ryu, and J. Baik, ``HASPO: Harmony search-based parameter optimization for just-in-time software defect prediction in maritime software,'' Appl. Sci., vol. 11, no. 5, Art. no. 2002, 2021, doi: \mbox{10.3390/app11052002}.
\bibitem{Mostafa2025} S. Mostafa, S. T. Cynthia, B. Roy, and D. Mondal, ``Feature transformation for improved software bug detection and commit classification,'' J. Syst. Softw., vol. 219, Art. no. 112205, Jan. 2025, doi: \mbox{10.1016/j.jss.2024.112205}.
\bibitem{Chen2025} F. Chen, Y. Ke, X. Liu, and Q. Wei, ``Multi-project just-in-time software defect prediction based on multi-task learning for mobile applications,'' in Proc. IEEE Conf. Softw. Testing, Verification Validation (ICST), 2025, pp. 127--137, doi: \mbox{10.1109/ICST62969.2025.10989004}.
\bibitem{Shehab2023} M. A. Shehab, A. Hamou-Lhadj, and V. S. P. Gunda, ``JITBoost: Boosting just-in-time defect prediction using Boolean combination of classifiers,'' in Proc. IEEE 23rd Int. Conf. Softw. Qual., Rel. Secur. (QRS), 2023, pp. 95--104, doi: \mbox{10.1109/QRS60937.2023.00019}.
\bibitem{Chitsazian2026} Z. Chitsazian, S. Sedighian Kashi, and A. Nikanjam, ``Detecting concept drift in just-in-time software defect prediction using model interpretation,'' Int. J. Syst. Assur. Eng. Manag., vol. 17, no. 9, pp. 2701--2722, 2026, doi: \mbox{10.1007/s13198-025-03122-7}.
\bibitem{Patel2024} H. J. Patel, B. Adams, and A. E. Hassan, ``Post deployment recycling of machine learning models,'' Empir. Softw. Eng., vol. 29, no. 4, Art. no. 100, 2024, doi: \mbox{10.1007/s10664-024-10492-2}.
\bibitem{Tian2026} D. Tian, L. Song, and X. Su, ``Just-in-time quality assurance with assistance of automated defect prediction and code testing,'' ACM Trans. Softw. Eng. Methodol., vol. 35, no. 9, pp. 1--46, 2026, doi: \mbox{10.1145/3779653}.
\bibitem{Song2023} L. Song and L. L. Minku, ``A procedure to continuously evaluate predictive performance of just-in-time software defect prediction models during software development,'' IEEE Trans. Softw. Eng., vol. 49, no. 2, pp. 646--666, Feb. 2023, doi: \mbox{10.1109/TSE.2022.3158831}.
\bibitem{Veras2025} R. C. Veras, G. G. Cabral, and A. L. I. de Oliveira, ``Semantic SZZ: Mitigating the impact of misclassified corrective changes in just-in-time software defect prediction,'' in Proc. IEEE Int. Conf. Syst., Man, Cybern. (SMC), 2025, pp. 769--774, doi: \mbox{10.1109/SMC58881.2025.11342537}.
\end{thebibliography}


\end{document}
